\paragraph{Weighted monic-cubic first integrals.}
The next result treats polynomial sections of a cubic first integral
whose denominator depends only on the evaluation coordinate.  It bounds
lists for arbitrary received words, including words supported on its
critical fibers.

\begin{theorem}\label{thm:weighted-cubic-first-integral-list}
Let $\mathbb K$ have characteristic zero or characteristic
$p>\max\{3,2D\}$, where $0\le D<n$.  Let $0\ne H\in\mathbb K[X]$ and
\[
 F(X,u)=u^3+a_2(X)u^2+a_1(X)u+a_0(X),
 \qquad \deg a_i\le(3-i)D\quad(0\le i\le2).
\]
Consider the section family
\[
 \mathcal B=\{P\in\mathbb K[X]:\deg P\le D,\quad
                   F(X,P)=cH\text{ for some }c\in\mathbb K\}.
\]
On any $n$ distinct evaluation coordinates, every received word has at most
\[
 \left\lfloor9+\frac{3+108D/n}{\eta}\right\rfloor
\]
members of $\mathcal B$ agreeing on at least $D+\eta n$ coordinates,
for every $\eta>0$.  Coordinates and received symbols may lie in an
extension field.
\end{theorem}

This concerns the displayed first-integral family, not arbitrary cubic
differential equations or all Reed--Solomon codewords.  In particular,
the characteristic and weighted coefficient hypotheses must both be
checked; this theorem does not cover the Dickson differential-equation
bank merely because its equation has cubic jet degree.

\begin{proof}
Extend constants to an algebraically closed field $k$ containing all
coordinates and received symbols, and put $K=k(X)$.  Each section
has a unique label $c$, and each label has at most three sections.
The case $D=0$ follows because distinct constants have disjoint agreement
sets.  Assume $D\ge1$.  If no nonzero label occurs, there are at most
three sections.  Otherwise, writing $N=\deg H$, the weighted degree
caps imply $N\le3D$.

\emph{The critical scheme.}
Set
\[
 B=HF_X-H'F,\qquad I=(F_u,B)\subset k[X,u].
\]
First suppose $I$ is zero-dimensional, allowing the unit ideal, and
write $\ell=\dim_k k[X,u]/I$.  The bidegrees in $(X,u)$ are at most
$(2D,2)$ and $(N+3D-1,3)$, respectively.  Homogenize each nonzero
generator in $\mathbb P^1\times\mathbb P^1$ using its own actual
bidegrees.  Neither closure contains either boundary divisor: such a
component would mean that the corresponding leading coefficient
polynomial was identically zero.  A common curve would therefore meet
the affine chart, contradicting zero dimensionality.  The intersection
product consequently bounds the affine scheme length by
\begin{equation}\label{eq:cubic-critical-length}
 \ell\le (2D)3+2(N+3D-1)=2N+12D-2.
\end{equation}
The unit-ideal case is immediate.  All arguments use the original ideal
$I$, without removing common factors from its generators.

For a section $P$ put $A_P=F_u(X,P)$ and $b_P=3P+a_2$.  Differentiating
$F(X,P)=cH$ gives
\[
 B(X,P)=-H A_P P',\qquad
 k[X,u]/(I,u-P)\simeq k[X]/(A_P).
\]
In particular $A_P\ne0$, since otherwise the entire section graph
would lie in the critical scheme.  At every $x$, the order of $A_P$
is at most the local length of this scheme at $(x,P(x))$.

\emph{Two valuation estimates.}
Taylor expansion, with no division by factorials, reads
\[
 F(X,P+Z)=cH+A_PZ+b_PZ^2+Z^3.
\]
Choose roots $\delta$ of $A_P+2b_P\delta+3\delta^2=0$.  They give
the two fixed critical roots $P+\delta$ of $F_u$, and
\begin{equation}\label{eq:cubic-critical-correction}
 F(X,P+\delta)-cH=-b_P\delta^2-2\delta^3.
\end{equation}
All valuations below are taken in one fixed valued algebraic closure
at the specified place.  Only quadratic valuation comparisons are
used; no positive-characteristic Puiseux expansion is needed.

At infinity, use degree equal to minus the valuation and put
$d=\deg A_P$, $e=\deg b_P$, with $e=-\infty$ when $b_P=0$.
If $2e>d$, one can choose $\deg\delta=d-e$, and the right side of
\eqref{eq:cubic-critical-correction} has degree at most
$2d-e<3d/2$.  If $2e\le d$, both roots have degree at most $d/2$,
so the correction has degree at most $3d/2$.  Thus, when $d<2N/3$,
one of the two fixed critical values $F(X,r)/H$ is regular at infinity
with residue $c$.  There are at most two such residues.  Apart from
at most two labels, therefore,
\begin{equation}\label{eq:cubic-infinity-cost}
 \deg A_P\ge\lceil2N/3\rceil.
\end{equation}

At a root $x$ of $H$ of multiplicity $m$, put
$d=\operatorname{ord}_x A_P$ and $e=\operatorname{ord}_x b_P$.
If $2e<d$, the root of larger order has order $d-e$, and the correction
in \eqref{eq:cubic-critical-correction} has order at least
$2d-e>3d/2$.  If $2e\ge d$, choose a root of order at least $d/2$;
the correction has order at least $3d/2$.  Hence
$d>\lfloor2m/3\rfloor$ forces a critical value $F(X,r)/H$ to have
residue $c$ at $x$.  For each fixed $x$, this can occur for at most
two labels.  Fractional valuations and repeated critical roots cause
no change to either estimate.

\emph{Charging multiplicities at repeated roots of $H$.}
Let $n_0$ be the number of distinct roots of $H$, and suppose $n_0>D$.
Choose one section for each nonzero label not excluded by
\eqref{eq:cubic-infinity-cost}.  Charge it
\[
 C(P)=\sum_{H(x)\ne0}\operatorname{ord}_x A_P+
 \sum_{H(x)=0}\max\{0,\operatorname{ord}_x A_P
                     -\lfloor2\operatorname{ord}_x H/3\rfloor\}.
\]
Off $H=0$, a critical point fixes its label by $c=F/H$, so different
labels have disjoint charged support.  Above each root of $H$, only
two labels have positive excess.  The local quotient computation
therefore gives $\sum_P C(P)\le2\ell$, even when both sections pass
through the same critical point.  Moreover,
\begin{align*}
 C(P)&\ge\lceil2N/3\rceil
             -\sum_{H(x)=0}\lfloor2\operatorname{ord}_x H/3\rfloor\\
     &\ge n_0-\lfloor N/3\rfloor\ge n_0-D>0,
\end{align*}
using $\lfloor2m/3\rfloor\le m-1$ for $m\ge1$.
After adding the two exceptional labels and the zero label, and then
multiplying the number of labels by three, we obtain the total-bank bound
\begin{equation}\label{eq:cubic-root-bank-bound}
 |\mathcal B|\le
 9+3\left\lfloor\frac{4N+24D-4}{n_0-\lfloor N/3\rfloor}\right\rfloor
 \le9+\frac{108D}{n_0-D}.
\end{equation}

\emph{Positive-dimensional critical loci.}
We next treat the case excluded above.  Since $F_u$ has constant leading
$u$-coefficient $3$, a positive-dimensional common component cannot be
vertical.  A root $r$ of such a component satisfies
\[
 F_u(X,r)=0,\qquad (F(X,r)/H)'=0.
\]
We use the elementary height fact that a rational function of height
less than $p$ with zero derivative is constant: it belongs to $K^p$,
and a nonconstant $p$-th power has height at least $p$.
We also need a distinct-root budget.  If $A/H$ is a nonconstant
$p$-th power with $\deg A,\deg H\le3D$, write
\[
 A=G a^p,\qquad H=G b^p,\qquad \gcd(a,b)=1,
 \qquad M=\max\{\deg a,\deg b\}\ge1.
\]
Then $\deg G+pM\le3D$, and the number $n_0$ of distinct roots of $H$
satisfies
\[
 n_0\le\deg G+\deg b\le3D-(p-1)M\le3D-p+1\le D.
\]
In this case a received symbol off $H=0$ fixes a label and matches
at most three candidates.  A list of $L$ candidates with at least
$a\ge D+\eta n$ agreements thus satisfies $La\le LD+3n$, and hence
$L\le3n/(a-D)\le3/\eta$.  We may finish immediately whenever
such a nonconstant Frobenius ratio occurs.

If $r\in K$, monic integrality makes $r$ polynomial, and the weighted
caps give $\deg r\le D$, so $\deg F(X,r)\le3D$.  In positive
characteristic a nonconstant ratio $F(X,r)/H$ is covered by the
distinct-root budget just proved.  Otherwise the ratio is constant,
say $c_0$, as it always is in characteristic zero.
The repeated root of the cubic $F-c_0H$ then gives
\begin{equation}\label{eq:cubic-repeated-fiber}
 F-c_0H=(u-R)^2(u-S),\qquad R,S\in k[X],\quad
 \deg R,\deg S\le D.
\end{equation}

For completeness, an irreducible quadratic critical component cannot
produce an additional large family under our sharper characteristic
bound.  Depress the cubic, writing
\[
 F=z^3-3gz+b,\qquad z=u+a_2/3,\qquad
 \deg g\le2D,\quad\deg b\le3D.
\]
Here $g$ is a nonsquare in $K$.  Both conjugate critical values have
zero derivative.  Their trace gives $(b/H)'=0$.  A nonconstant ratio
in positive characteristic is again covered by the distinct-root
budget, since $\deg b\le3D$.  In all remaining cases $b=\beta H$
for a constant $\beta$.
Their difference gives $3g'H-2gH'=0$.
A section with $\lambda=c-\beta\ne0$ would satisfy
$z(z^2-3g)=\lambda H$.  Differentiating and eliminating $H'$ yields
\[
 (z^2-g)(2gz'-g'z)=0.
\]
The first factor is nonzero because $g$ is a nonsquare, and $z\ne0$.
Thus $(g/z^2)'=0$.  Its height is at most $2D<p$, so $g/z^2$ is
constant, contradicting nonsquareness over the algebraically closed
constant field.  All sections therefore have label $\beta$, giving
at most three.  The same argument in characteristic zero uses the
usual constant field of $K$ instead of a height bound.

It remains to analyze \eqref{eq:cubic-repeated-fiber}.  If $R=S$, any
section outside this fiber makes $H$ a constant times a polynomial cube,
and all sections lie in an affine family $R+t h$ with $\deg h\le D$.
If $R\ne S$, put $h=R-S$, $Y=(P-R)/h$, and $q=H/h^3$.  Then
\begin{equation}\label{eq:cubic-three-cover-equation}
 Y^2(Y+1)=(c-c_0)q.
\end{equation}
If $q$ is constant, $Y$ is constant and again all sections lie in
$R+t h$.  In an affine family of degree-at-most-$D$ polynomials,
different members can match the same received symbol only at zeros
of $h$, of which there are at most $D$.  A list of $L$ members with
at least $a>D$ agreements therefore satisfies $L(a-D)\le n$.

If $q$ is nonconstant, there are at most two distinct nonzero labels
$\lambda=c-c_0$ for which \eqref{eq:cubic-three-cover-equation} has a
rational solution.  Here is a self-contained cover argument.  Three
such labels would give three degree-three covers of the $t$-line
\[
 Y_i^2(Y_i+1)=\lambda_i t.
\]
Each has branch points $0,4/(27\lambda_i),\infty$, with inertia
types a transposition, a transposition, and a three-cycle.  Its Galois
closure has group $S_3$.  The compositum group projects onto all three
$S_3$ factors.  At $4/(27\lambda_i)$ inertia is a transposition
supported only in factor $i$, since the other covers are unramified
there.  Conjugation supplies that entire $S_3$ factor, proving that
the compositum group is $S_3^3$.  The original fiber product is thus
connected of degree $27$.

All ramification is tame in characteristic greater than three.  On the
$27$ sheets the inertia at zero has one fixed point and thirteen
two-cycles, so index $13$; at infinity it has nine three-cycles, so
index $18$; and each separate finite nonzero branch point has index
$9$.  The smooth normalization $C$ consequently has
\[
 2g(C)-2=-54+13+18+3\cdot9=4,
 \qquad g(C)=3.
\]
The proposed rational solutions would embed $k(C)$ in $k(X)$ via
$t=q(X)$, contradicting L\"uroth's theorem.  This also rules out
inseparable parametrizations.  There are therefore at most six
sections outside the repeated fiber and at most two within it, for a
total of eight.  This completes the positive-dimensional case.

\emph{From the total-bank bound to arbitrary received words.}
Only the zero-dimensional case remains.  Let $L$ selected candidates
have at least $a\ge D+\eta n$ agreements.  Let $r$ be the number of
evaluation coordinates where $H$ vanishes, so $r\le n_0$.  At any
other coordinate $x$, the received value $w(x)$ determines the label
$F(x,w(x))/H(x)$, and hence matches at most three candidates.  Thus
\[
 La\le Lr+3(n-r)\le Ln_0+3n,
 \qquad n_0-D\ge\eta n-3n/L.
\]
If $L\le\max\{9,3/\eta\}$ the asserted bound is immediate.
Otherwise the last quantity is positive, and
\eqref{eq:cubic-root-bank-bound} implies
\[
 L\le9+\frac{108D}{\eta n-3n/L}.
\]
Multiplication by the positive denominator gives
\[
 \eta L\le9\eta+3+108D/n-27/L
            <9\eta+3+108D/n.
\]
Taking the stated floor proves the result.  All preceding small-bank
and affine-family cases satisfy the same bound.
\end{proof}

\paragraph{A distinction for the Dickson family.}
For the complete nonzero Dickson bank, $n=p-1$, $D=n/4-1$, and
$n/2$ candidates agree with the displayed received word on $3n/8$
coordinates.  If this bank lay in a common first-integral family as
above, taking $\eta=1/8+1/n$ would give
\[
 \frac n2\le\left\lfloor249-\frac{2784}{n+8}\right\rfloor.
\]
This fails for $n\ge488$, hence for every admissible prime
$p\equiv1\pmod8$ with $p\ge521$.  Thus its primitive cubic differential
equation does not arise from a weighted monic-cubic first integral of
this form.  This does not exclude higher-degree first integrals or
rational denominators depending on the candidate value.
